\section{Conclusion}

In this work, we studied multi-contrast MRI reconstruction from undersampled k-space data. We first established a single-modal U-Net baseline for Task 2, which effectively suppresses aliasing artifacts but tends to produce over-smoothed reconstructions. To overcome this limitation, we proposed a Multi-modal Unrolled Reconstruction Network for Task 3, which leverages the fully sampled T1 modality as anatomical guidance for undersampled T2 reconstruction. By combining cross-attention-based feature fusion, iterative unrolled refinement, and explicit k-space data consistency, the proposed method achieves more faithful and perceptually realistic reconstructions.

Extensive experiments show that our Task 3 model consistently outperforms the Task 2 baseline across PSNR, SSIM, LPIPS, and DISTS. Qualitative results further demonstrate that the proposed framework better preserves anatomical boundaries, tumor regions, and fine structural details while reducing residual undersampling artifacts. In addition, our analysis of PSNR reveals its limitation in reflecting perceptual image quality, motivating the use of perceptual and structure-aware metrics. The ablation study confirms the effectiveness of the proposed components, including wavelet loss, perceptual loss, and cross-attention-based multi-modal fusion.

Despite these improvements, challenging pathological cases still exhibit mild over-smoothing and boundary artifacts. Future work may explore stronger edge-aware constraints, more robust multi-modal registration, and advanced frequency-domain modeling. Overall, our results demonstrate that integrating multi-modal anatomical priors with unrolled data consistency and perceptually motivated objectives is an effective direction for high-quality accelerated MRI reconstruction.